\begin{figure}[htbp]
\centering
\begin{tikzpicture}[
    >=Latex,
    font=\small,
    box/.style={draw, semithick, rounded corners=2pt, align=center,
                minimum height=10mm, minimum width=34mm, inner sep=4pt, fill=black!2},
    key/.style={box, fill=black!7},
    small/.style={draw, semithick, rounded corners=2pt, align=center,
                  minimum height=9mm, minimum width=27mm, inner sep=3pt, fill=black!2},
    arr/.style={->, semithick}
]
\node[font=\bfseries\normalsize] at (-4.1,3.25) {Generative};
\node[font=\bfseries\normalsize] at (4.1,3.25) {Contrastive};
\node[key] (gx) at (-4.1,2.35) {Input $\mathbf{x}$};
\node[box] (ge) at (-4.1,0.95) {Encoder};
\node[box] (gz) at (-4.1,-0.45) {Latent $\mathbf{z}$};
\node[box] (gd) at (-4.1,-1.85) {Decoder};
\node[key] (gr) at (-4.1,-3.25) {Reconstruction $\hat{\mathbf{x}}$};
\draw[arr] (gx)--(ge); \draw[arr] (ge)--(gz); \draw[arr] (gz)--(gd); \draw[arr] (gd)--(gr);

\node[key] (cx) at (4.1,2.35) {Input $\mathbf{x}$};
\node[small] (ca) at (2.35,0.85) {Augmented view A};
\node[small] (cb) at (5.85,0.85) {Augmented view B};
\node[small] (cea) at (2.35,-0.55) {Shared encoder};
\node[small] (ceb) at (5.85,-0.55) {Shared encoder};
\node[small] (cza) at (2.35,-1.95) {$\mathbf{z}_a$};
\node[small] (czb) at (5.85,-1.95) {$\mathbf{z}_b$};
\node[key] (closs) at (4.1,-3.45) {Contrastive objective};
\draw[arr] (cx)--(ca); \draw[arr] (cx)--(cb);
\draw[arr] (ca)--(cea); \draw[arr] (cb)--(ceb);
\draw[arr] (cea)--(cza); \draw[arr] (ceb)--(czb);
\draw[<->, semithick] (cza) -- node[midway, above=10pt, fill=white, inner sep=1pt]{positive pair} (czb);
\draw[arr] (cza)--(closs); \draw[arr] (czb)--(closs);
\draw[gray, dashed] (0,2.95)--(0,-4.05);
\end{tikzpicture}

\caption{Conceptual comparison between generative and contrastive representation learning. Generative models optimize reconstruction through a latent variable, whereas contrastive methods organize representations through similarity relationships between augmented views. Author's elaboration based on \cite{kingma2013auto,chen2020simple}.}
\label{fig:generative_contrastive_topologies}
\end{figure}
